\clearpage\fancyhead[L]{\small COMET.}\hypertarget{native-369-lower}{}\label{native-369-lower}
\begin{center}\textbf{COMET.}\par\vspace{4pt}\hrule\vspace{4pt}\fontsize{10}{12}\selectfont Monatliche Correspondenz. Band XXVII. Seite 388\;..\;389. 1813 April.\par\vspace{4pt}\hrule\vspace{4pt}\end{center}

\newcommand{\CometThreeSixNineA}{\hbox{\vrule width0pt height8.041667bp depth0pt\special{pdf:bcontent}\special{pdf:code q 0 g 8.583333 7.500000 m 8.083333 7.500000 l 8.083333 7.375000 l 7.916667 7.333333 l 7.916667 7.208333 l 7.750000 7.166667 l 7.750000 7.000000 l 7.541667 6.958333 l 7.541667 6.833333 l 7.375000 6.791667 l 7.333333 6.625000 l 7.000000 6.625000 l 7.000000 6.458333 l 6.833333 6.416667 l 6.833333 6.291667 l 6.666667 6.250000 l 6.625000 6.083333 l 6.458333 6.083333 l 6.458333 5.916667 l 6.291667 5.875000 l 6.250000 5.708333 l 6.125000 5.708333 l 6.125000 5.375000 l 5.958333 5.333333 l 5.916667 5.166667 l 5.750000 5.166667 l 5.750000 5.041667 l 5.583333 5.000000 l 5.541667 4.833333 l 5.208333 4.833333 l 5.208333 4.666667 l 5.041667 4.625000 l 5.041667 4.500000 l 4.875000 4.458333 l 4.833333 4.291667 l 4.500000 4.291667 l 4.500000 4.125000 l 4.333333 4.083333 l 4.291667 3.916667 l 4.125000 3.916667 l 4.125000 3.791667 l 3.958333 3.750000 l 3.916667 3.541667 l 3.791667 3.541667 l 3.750000 3.750000 l 3.541667 3.791667 l 3.541667 4.083333 l 3.750000 4.125000 l 3.791667 4.333333 l 4.083333 4.333333 l 4.083333 4.458333 l 4.291667 4.500000 l 4.291667 4.625000 l 4.458333 4.666667 l 4.458333 4.833333 l 4.625000 4.875000 l 4.666667 5.041667 l 4.833333 5.041667 l 4.875000 5.208333 l 5.166667 5.208333 l 5.208333 5.375000 l 5.375000 5.375000 l 5.375000 5.541667 l 5.541667 5.583333 l 5.541667 5.708333 l 5.708333 5.750000 l 5.750000 5.916667 l 5.916667 5.916667 l 5.916667 6.083333 l 6.083333 6.125000 l 6.083333 6.250000 l 6.250000 6.291667 l 6.291667 6.458333 l 6.625000 6.458333 l 6.625000 6.625000 l 6.791667 6.666667 l 6.791667 6.791667 l 6.958333 6.833333 l 7.000000 7.000000 l 7.166667 7.000000 l 7.208333 7.208333 l 7.500000 7.208333 l 7.541667 7.375000 l 7.708333 7.375000 l 7.708333 7.708333 l 7.541667 7.708333 l 7.500000 7.875000 l 6.666667 7.875000 l 6.666667 7.750000 l 6.458333 7.708333 l 6.458333 7.541667 l 6.291667 7.500000 l 6.291667 7.375000 l 6.125000 7.333333 l 6.083333 7.166667 l 5.958333 7.166667 l 5.958333 7.000000 l 5.750000 6.958333 l 5.750000 6.833333 l 5.583333 6.791667 l 5.541667 6.625000 l 5.416667 6.625000 l 5.416667 6.458333 l 5.208333 6.416667 l 5.208333 6.291667 l 5.041667 6.250000 l 5.000000 6.083333 l 4.875000 6.083333 l 4.875000 5.916667 l 4.666667 5.875000 l 4.666667 5.750000 l 4.500000 5.708333 l 4.458333 5.541667 l 4.333333 5.541667 l 4.291667 5.333333 l 4.125000 5.333333 l 4.125000 5.041667 l 3.958333 5.000000 l 3.916667 4.833333 l 3.583333 4.833333 l 3.583333 4.666667 l 3.416667 4.625000 l 3.375000 4.458333 l 3.041667 4.458333 l 3.041667 4.333333 l 2.875000 4.291667 l 2.833333 4.333333 l 2.833333 4.625000 l 3.000000 4.666667 l 3.000000 4.833333 l 3.208333 4.875000 l 3.250000 5.041667 l 3.541667 5.041667 l 3.583333 5.208333 l 3.750000 5.208333 l 3.750000 5.333333 l 3.916667 5.375000 l 3.916667 5.541667 l 4.083333 5.583333 l 4.125000 5.750000 l 4.291667 5.750000 l 4.291667 5.875000 l 4.458333 5.916667 l 4.458333 6.083333 l 4.625000 6.125000 l 4.666667 6.291667 l 4.833333 6.291667 l 4.833333 6.416667 l 5.000000 6.458333 l 5.000000 6.625000 l 5.166667 6.666667 l 5.208333 6.833333 l 5.375000 6.833333 l 5.375000 6.958333 l 5.541667 7.000000 l 5.541667 7.166667 l 5.708333 7.208333 l 5.708333 7.500000 l 5.916667 7.541667 l 5.916667 7.875000 l 5.750000 7.875000 l 5.708333 8.041667 l 5.041667 8.041667 l 5.000000 7.875000 l 4.666667 7.875000 l 4.666667 7.750000 l 4.500000 7.708333 l 4.500000 7.375000 l 4.333333 7.333333 l 4.291667 7.166667 l 4.125000 7.166667 l 4.125000 7.000000 l 3.958333 6.958333 l 3.958333 6.666667 l 3.791667 6.625000 l 3.750000 6.416667 l 3.583333 6.416667 l 3.583333 6.125000 l 3.416667 6.083333 l 3.416667 5.916667 l 3.250000 5.875000 l 3.208333 5.708333 l 3.041667 5.708333 l 3.041667 5.583333 l 2.875000 5.541667 l 2.833333 5.333333 l 2.708333 5.333333 l 2.708333 5.041667 l 2.500000 5.000000 l 2.500000 4.875000 l 2.333333 4.833333 l 2.333333 4.666667 l 2.166667 4.625000 l 2.125000 4.458333 l 1.791667 4.458333 l 1.750000 4.291667 l 1.250000 4.291667 l 1.208333 4.083333 l 0.916667 4.083333 l 0.916667 3.958333 l 0.708333 3.916667 l 0.708333 3.791667 l 0.541667 3.750000 l 0.541667 3.583333 l 0.375000 3.541667 l 0.375000 3.041667 l 0.166667 3.000000 l 0.166667 1.958333 l 0.000000 1.916667 l 0.000000 1.416667 l 0.166667 1.375000 l 0.166667 1.083333 l 0.333333 1.083333 l 0.375000 0.875000 l 0.500000 0.875000 l 0.541667 0.708333 l 0.708333 0.666667 l 0.708333 0.541667 l 0.875000 0.541667 l 0.916667 0.375000 l 1.208333 0.375000 l 1.250000 0.166667 l 1.750000 0.166667 l 1.791667 0.000000 l 3.208333 0.000000 l 3.208333 0.125000 l 3.375000 0.166667 l 3.375000 0.333333 l 3.541667 0.375000 l 3.583333 0.541667 l 3.750000 0.541667 l 3.750000 0.666667 l 3.916667 0.708333 l 3.958333 0.875000 l 4.083333 0.875000 l 4.083333 1.375000 l 4.291667 1.416667 l 4.291667 2.666667 l 4.083333 2.708333 l 4.083333 3.208333 l 4.125000 3.250000 l 4.458333 3.250000 l 4.500000 3.416667 l 4.833333 3.416667 l 4.833333 3.541667 l 5.000000 3.583333 l 5.041667 3.791667 l 5.375000 3.791667 l 5.375000 3.916667 l 5.541667 3.958333 l 5.583333 4.125000 l 5.916667 4.125000 l 5.916667 4.291667 l 6.083333 4.333333 l 6.125000 4.500000 l 6.416667 4.500000 l 6.458333 4.666667 l 6.625000 4.666667 l 6.625000 4.833333 l 6.791667 4.875000 l 6.791667 5.000000 l 6.958333 5.041667 l 7.000000 5.208333 l 7.166667 5.208333 l 7.166667 5.333333 l 7.333333 5.375000 l 7.333333 5.541667 l 7.500000 5.583333 l 7.541667 5.750000 l 7.875000 5.750000 l 7.875000 5.875000 l 8.041667 5.916667 l 8.041667 6.083333 l 8.250000 6.125000 l 8.250000 6.250000 l 8.416667 6.291667 l 8.458333 6.458333 l 8.583333 6.458333 l h 1.625000 3.583333 m 2.666667 3.583333 l 2.708333 3.416667 l 3.000000 3.416667 l 3.041667 3.250000 l 3.208333 3.250000 l 3.250000 3.041667 l 3.416667 3.000000 l 3.416667 2.333333 l 3.250000 2.291667 l 3.250000 1.625000 l 3.041667 1.583333 l 3.041667 1.416667 l 2.875000 1.375000 l 2.875000 1.250000 l 2.708333 1.208333 l 2.666667 1.041667 l 2.500000 1.041667 l 2.458333 0.833333 l 1.625000 0.833333 l 1.583333 1.041667 l 1.458333 1.041667 l 1.416667 1.208333 l 1.250000 1.208333 l 1.208333 1.375000 l 1.041667 1.416667 l 1.041667 1.583333 l 0.875000 1.625000 l 0.875000 1.750000 l 0.666667 1.791667 l 0.666667 2.458333 l 0.875000 2.500000 l 0.875000 2.833333 l 1.041667 2.875000 l 1.041667 3.000000 l 1.208333 3.041667 l 1.208333 3.208333 l 1.416667 3.250000 l 1.416667 3.375000 l 1.583333 3.416667 l h f* Q}\special{pdf:econtent}\kern8.583333bp}}
\newcommand{\CometThreeSixNineB}{\hbox{\vrule width0pt height8.083333bp depth0pt\special{pdf:bcontent}\special{pdf:code q 0 g 8.416667 7.541667 m 8.083333 7.541667 l 8.083333 7.375000 l 7.916667 7.333333 l 7.916667 7.208333 l 7.750000 7.166667 l 7.708333 7.000000 l 7.583333 7.000000 l 7.541667 6.791667 l 7.375000 6.791667 l 7.333333 6.625000 l 7.041667 6.625000 l 7.000000 6.666667 l 7.000000 7.000000 l 7.166667 7.041667 l 7.208333 7.208333 l 7.541667 7.208333 l 7.541667 7.333333 l 7.708333 7.375000 l 7.708333 7.708333 l 7.583333 7.708333 l 7.541667 7.875000 l 6.666667 7.875000 l 6.625000 7.708333 l 6.291667 7.708333 l 6.250000 7.541667 l 5.958333 7.541667 l 5.916667 7.708333 l 5.750000 7.708333 l 5.708333 7.875000 l 5.541667 7.916667 l 5.541667 8.083333 l 5.208333 8.083333 l 5.166667 7.875000 l 4.875000 7.875000 l 4.833333 7.708333 l 4.666667 7.708333 l 4.666667 7.583333 l 4.500000 7.541667 l 4.500000 7.375000 l 4.333333 7.333333 l 4.291667 7.166667 l 4.125000 7.166667 l 4.125000 7.041667 l 3.958333 7.000000 l 3.958333 6.833333 l 3.791667 6.791667 l 3.750000 6.625000 l 3.583333 6.625000 l 3.583333 6.291667 l 3.416667 6.250000 l 3.416667 6.125000 l 3.250000 6.083333 l 3.208333 5.916667 l 3.041667 5.916667 l 3.041667 5.750000 l 2.875000 5.708333 l 2.875000 5.416667 l 2.708333 5.375000 l 2.708333 5.208333 l 2.541667 5.166667 l 2.500000 5.000000 l 2.333333 5.000000 l 2.333333 4.875000 l 2.166667 4.833333 l 2.125000 4.666667 l 2.000000 4.666667 l 1.958333 4.458333 l 1.791667 4.458333 l 1.750000 4.291667 l 1.250000 4.291667 l 1.208333 4.125000 l 0.916667 4.125000 l 0.916667 3.958333 l 0.708333 3.916667 l 0.708333 3.791667 l 0.541667 3.750000 l 0.500000 3.583333 l 0.375000 3.583333 l 0.375000 3.250000 l 0.166667 3.208333 l 0.166667 2.708333 l 0.000000 2.666667 l 0.000000 1.625000 l 0.166667 1.583333 l 0.166667 1.250000 l 0.333333 1.250000 l 0.375000 1.083333 l 0.541667 1.041667 l 0.541667 0.708333 l 0.666667 0.708333 l 0.708333 0.541667 l 1.041667 0.541667 l 1.083333 0.375000 l 1.250000 0.333333 l 1.250000 0.208333 l 1.458333 0.166667 l 1.458333 0.000000 l 3.000000 0.000000 l 3.000000 0.166667 l 3.208333 0.208333 l 3.208333 0.333333 l 3.375000 0.375000 l 3.375000 0.500000 l 3.541667 0.541667 l 3.583333 0.708333 l 3.750000 0.708333 l 3.750000 0.875000 l 3.916667 0.916667 l 3.958333 1.083333 l 4.083333 1.083333 l 4.083333 1.416667 l 4.291667 1.458333 l 4.291667 2.500000 l 4.083333 2.541667 l 4.083333 3.041667 l 4.291667 3.083333 l 4.333333 3.250000 l 4.625000 3.250000 l 4.666667 3.416667 l 5.166667 3.416667 l 5.166667 3.583333 l 5.375000 3.625000 l 5.416667 3.791667 l 5.708333 3.791667 l 5.708333 3.916667 l 5.916667 3.958333 l 5.916667 4.125000 l 6.083333 4.166667 l 6.083333 4.291667 l 6.250000 4.333333 l 6.291667 4.500000 l 6.458333 4.500000 l 6.458333 4.666667 l 6.625000 4.708333 l 6.625000 4.833333 l 6.791667 4.875000 l 6.833333 5.041667 l 7.000000 5.041667 l 7.000000 5.166667 l 7.166667 5.208333 l 7.166667 5.375000 l 7.333333 5.416667 l 7.375000 5.583333 l 7.541667 5.583333 l 7.541667 5.708333 l 7.708333 5.750000 l 7.750000 5.958333 l 8.041667 5.958333 l 8.083333 6.125000 l 8.250000 6.125000 l 8.250000 6.250000 l 8.416667 6.291667 l 8.416667 6.458333 l 8.583333 6.500000 l 8.583333 6.791667 l 8.791667 6.833333 l 8.791667 7.166667 l 8.625000 7.166667 l 8.583333 7.333333 l 8.416667 7.375000 l h 1.458333 3.625000 m 2.125000 3.625000 l 2.166667 3.416667 l 2.833333 3.416667 l 2.875000 3.250000 l 3.000000 3.250000 l 3.041667 3.083333 l 3.208333 3.083333 l 3.250000 2.875000 l 3.416667 2.833333 l 3.416667 2.541667 l 3.250000 2.500000 l 3.250000 1.625000 l 3.041667 1.583333 l 3.041667 1.458333 l 2.875000 1.416667 l 2.875000 1.250000 l 2.708333 1.208333 l 2.666667 1.041667 l 2.166667 1.041667 l 2.125000 0.875000 l 1.250000 0.875000 l 1.208333 1.041667 l 1.041667 1.083333 l 1.041667 1.416667 l 0.916667 1.416667 l 0.875000 1.583333 l 0.666667 1.625000 l 0.666667 1.958333 l 0.500000 2.000000 l 0.500000 2.500000 l 0.666667 2.541667 l 0.666667 2.833333 l 0.875000 2.875000 l 0.875000 3.041667 l 1.041667 3.083333 l 1.041667 3.208333 l 1.208333 3.250000 l 1.250000 3.416667 l 1.416667 3.416667 l h 3.916667 5.416667 m 3.958333 5.375000 l 3.958333 5.041667 l 3.791667 5.000000 l 3.750000 4.833333 l 3.583333 4.833333 l 3.541667 4.666667 l 3.250000 4.666667 l 3.250000 4.500000 l 3.041667 4.458333 l 3.000000 4.500000 l 3.000000 4.833333 l 3.208333 4.875000 l 3.208333 5.000000 l 3.375000 5.041667 l 3.416667 5.208333 l 3.750000 5.208333 l 3.750000 5.375000 l h 6.833333 6.458333 m 6.833333 6.291667 l 6.666667 6.250000 l 6.666667 6.125000 l 6.500000 6.083333 l 6.458333 5.916667 l 6.291667 5.916667 l 6.291667 5.750000 l 6.125000 5.708333 l 6.125000 5.583333 l 5.958333 5.541667 l 5.916667 5.375000 l 5.750000 5.375000 l 5.708333 5.166667 l 5.416667 5.166667 l 5.375000 5.000000 l 5.208333 5.000000 l 5.208333 4.875000 l 5.041667 4.833333 l 5.000000 4.666667 l 4.875000 4.666667 l 4.875000 4.500000 l 4.666667 4.458333 l 4.666667 4.333333 l 4.500000 4.291667 l 4.500000 4.166667 l 4.333333 4.125000 l 4.333333 3.958333 l 4.125000 3.916667 l 4.125000 3.791667 l 3.958333 3.750000 l 3.916667 3.583333 l 3.583333 3.583333 l 3.541667 3.750000 l 3.375000 3.791667 l 3.375000 3.916667 l 3.541667 3.958333 l 3.541667 4.125000 l 3.750000 4.166667 l 3.791667 4.333333 l 4.083333 4.333333 l 4.125000 4.500000 l 4.291667 4.500000 l 4.333333 4.708333 l 4.625000 4.708333 l 4.666667 4.875000 l 4.833333 4.875000 l 4.833333 5.000000 l 5.000000 5.041667 l 5.041667 5.208333 l 5.166667 5.208333 l 5.166667 5.375000 l 5.375000 5.416667 l 5.375000 5.541667 l 5.541667 5.583333 l 5.541667 5.708333 l 5.708333 5.750000 l 5.708333 5.916667 l 5.916667 5.958333 l 5.916667 6.083333 l 6.083333 6.125000 l 6.125000 6.291667 l 6.458333 6.291667 l 6.500000 6.500000 l 6.791667 6.500000 l h 5.541667 7.041667 m 5.583333 7.000000 l 5.583333 6.666667 l 5.416667 6.625000 l 5.375000 6.458333 l 5.208333 6.458333 l 5.208333 6.291667 l 5.041667 6.250000 l 5.041667 6.125000 l 4.875000 6.083333 l 4.833333 5.916667 l 4.666667 5.916667 l 4.666667 5.750000 l 4.500000 5.708333 l 4.500000 5.583333 l 4.333333 5.541667 l 4.291667 5.375000 l 3.958333 5.375000 l 3.916667 5.541667 l 4.083333 5.583333 l 4.083333 5.916667 l 4.291667 5.958333 l 4.333333 6.125000 l 4.625000 6.125000 l 4.625000 6.250000 l 4.833333 6.291667 l 4.833333 6.458333 l 5.000000 6.500000 l 5.000000 6.625000 l 5.166667 6.666667 l 5.208333 6.833333 l 5.375000 6.833333 l 5.416667 7.041667 l h f* Q}\special{pdf:econtent}\kern8.791667bp}}
\newcommand{\CometThreeSixNineC}{\hbox{\vrule width0pt height8.250000bp depth0pt\special{pdf:bcontent}\special{pdf:code q 0 g 8.625000 7.541667 m 8.458333 7.541667 l 8.416667 7.708333 l 8.125000 7.708333 l 8.125000 7.583333 l 7.916667 7.541667 l 7.875000 7.375000 l 7.583333 7.375000 l 7.583333 7.041667 l 7.375000 7.000000 l 7.333333 6.833333 l 7.041667 6.833333 l 7.041667 6.666667 l 6.833333 6.625000 l 6.833333 6.500000 l 6.666667 6.458333 l 6.625000 6.291667 l 6.500000 6.291667 l 6.500000 6.125000 l 6.291667 6.083333 l 6.291667 5.958333 l 6.125000 5.916667 l 6.083333 5.750000 l 5.958333 5.750000 l 5.958333 5.583333 l 5.750000 5.541667 l 5.750000 5.416667 l 5.583333 5.375000 l 5.541667 5.208333 l 5.416667 5.208333 l 5.416667 5.041667 l 5.250000 5.000000 l 5.208333 4.833333 l 5.041667 4.833333 l 5.000000 4.666667 l 4.875000 4.666667 l 4.833333 4.458333 l 4.500000 4.458333 l 4.500000 4.333333 l 4.333333 4.291667 l 4.291667 4.125000 l 4.166667 4.125000 l 4.125000 3.916667 l 3.958333 3.916667 l 3.916667 3.750000 l 3.625000 3.750000 l 3.583333 3.916667 l 3.375000 3.958333 l 3.416667 4.166667 l 3.583333 4.166667 l 3.625000 4.333333 l 3.916667 4.333333 l 3.958333 4.500000 l 4.125000 4.500000 l 4.125000 4.666667 l 4.291667 4.708333 l 4.291667 4.833333 l 4.458333 4.875000 l 4.500000 5.041667 l 4.666667 5.041667 l 4.666667 5.208333 l 4.833333 5.250000 l 4.875000 5.416667 l 5.208333 5.416667 l 5.208333 5.541667 l 5.375000 5.583333 l 5.375000 5.750000 l 5.541667 5.791667 l 5.541667 5.916667 l 5.708333 5.958333 l 5.750000 6.125000 l 5.916667 6.125000 l 5.916667 6.291667 l 6.083333 6.333333 l 6.083333 6.458333 l 6.250000 6.500000 l 6.291667 6.666667 l 6.458333 6.666667 l 6.500000 6.875000 l 6.791667 6.875000 l 6.833333 7.041667 l 7.000000 7.041667 l 7.000000 7.166667 l 7.166667 7.208333 l 7.166667 7.375000 l 7.333333 7.416667 l 7.375000 7.583333 l 7.708333 7.583333 l 7.708333 7.708333 l 7.541667 7.750000 l 7.541667 8.083333 l 6.666667 8.083333 l 6.625000 7.916667 l 6.291667 7.916667 l 6.291667 7.750000 l 6.125000 7.708333 l 6.083333 7.541667 l 5.916667 7.583333 l 5.916667 7.916667 l 5.750000 7.916667 l 5.708333 8.083333 l 5.541667 8.125000 l 5.541667 8.250000 l 5.041667 8.250000 l 5.000000 8.083333 l 4.708333 8.083333 l 4.666667 7.916667 l 4.500000 7.916667 l 4.500000 7.583333 l 4.333333 7.541667 l 4.333333 7.416667 l 4.166667 7.375000 l 4.166667 7.208333 l 3.958333 7.166667 l 3.958333 7.041667 l 3.791667 7.000000 l 3.750000 6.833333 l 3.625000 6.833333 l 3.625000 6.500000 l 3.416667 6.458333 l 3.416667 6.333333 l 3.250000 6.291667 l 3.250000 6.125000 l 3.083333 6.083333 l 3.041667 5.916667 l 2.875000 5.916667 l 2.833333 5.750000 l 2.708333 5.750000 l 2.708333 5.416667 l 2.541667 5.375000 l 2.500000 5.208333 l 2.333333 5.208333 l 2.333333 5.041667 l 2.166667 5.000000 l 2.166667 4.875000 l 2.000000 4.833333 l 1.958333 4.666667 l 1.625000 4.666667 l 1.583333 4.458333 l 1.083333 4.458333 l 1.083333 4.333333 l 0.916667 4.291667 l 0.875000 4.125000 l 0.708333 4.125000 l 0.708333 3.958333 l 0.541667 3.916667 l 0.500000 3.750000 l 0.375000 3.750000 l 0.375000 3.458333 l 0.208333 3.416667 l 0.208333 2.708333 l 0.000000 2.666667 l 0.000000 1.458333 l 0.166667 1.458333 l 0.208333 1.291667 l 0.375000 1.250000 l 0.375000 0.916667 l 0.500000 0.916667 l 0.541667 0.750000 l 0.708333 0.708333 l 0.708333 0.541667 l 0.916667 0.500000 l 0.916667 0.375000 l 1.416667 0.375000 l 1.458333 0.208333 l 1.958333 0.208333 l 2.000000 0.000000 l 2.833333 0.000000 l 2.833333 0.166667 l 3.041667 0.208333 l 3.083333 0.375000 l 3.375000 0.375000 l 3.416667 0.541667 l 3.583333 0.541667 l 3.583333 0.708333 l 3.750000 0.750000 l 3.750000 1.041667 l 3.916667 1.083333 l 3.958333 1.291667 l 4.125000 1.291667 l 4.125000 1.791667 l 4.291667 1.833333 l 4.291667 2.500000 l 4.125000 2.541667 l 4.125000 2.875000 l 3.958333 2.875000 l 3.916667 3.041667 l 4.125000 3.083333 l 4.125000 3.208333 l 4.291667 3.250000 l 4.291667 3.416667 l 4.458333 3.458333 l 4.500000 3.625000 l 4.666667 3.625000 l 4.708333 3.791667 l 5.000000 3.791667 l 5.041667 3.958333 l 5.375000 3.958333 l 5.375000 4.125000 l 5.541667 4.166667 l 5.541667 4.291667 l 5.708333 4.333333 l 5.750000 4.500000 l 5.916667 4.500000 l 5.916667 4.666667 l 6.083333 4.708333 l 6.125000 4.875000 l 6.458333 4.875000 l 6.458333 5.000000 l 6.625000 5.041667 l 6.625000 5.208333 l 6.791667 5.250000 l 6.833333 5.416667 l 7.166667 5.416667 l 7.166667 5.541667 l 7.333333 5.583333 l 7.333333 5.750000 l 7.541667 5.791667 l 7.541667 5.916667 l 7.708333 5.958333 l 7.750000 6.125000 l 7.875000 6.125000 l 7.916667 6.333333 l 8.250000 6.333333 l 8.250000 6.458333 l 8.416667 6.500000 l 8.458333 6.666667 l 8.625000 6.666667 l h 1.458333 3.791667 m 2.500000 3.791667 l 2.541667 3.625000 l 2.833333 3.625000 l 2.875000 3.458333 l 3.208333 3.458333 l 3.250000 3.416667 l 3.250000 3.083333 l 3.416667 3.041667 l 3.416667 2.541667 l 3.250000 2.500000 l 3.250000 1.833333 l 3.083333 1.791667 l 3.041667 1.583333 l 2.875000 1.583333 l 2.875000 1.291667 l 2.708333 1.250000 l 2.666667 1.041667 l 1.250000 1.041667 l 1.208333 1.250000 l 1.083333 1.250000 l 1.041667 1.416667 l 0.875000 1.458333 l 0.875000 1.791667 l 0.666667 1.833333 l 0.666667 3.041667 l 0.875000 3.083333 l 0.875000 3.208333 l 1.041667 3.250000 l 1.041667 3.416667 l 1.208333 3.458333 l 1.250000 3.625000 l 1.416667 3.625000 l h 5.750000 7.375000 m 5.750000 7.041667 l 5.583333 7.000000 l 5.541667 6.833333 l 5.416667 6.833333 l 5.416667 6.666667 l 5.250000 6.625000 l 5.208333 6.458333 l 5.041667 6.458333 l 5.041667 6.333333 l 4.875000 6.291667 l 4.875000 6.125000 l 4.708333 6.083333 l 4.666667 5.916667 l 4.500000 5.916667 l 4.500000 5.791667 l 4.333333 5.750000 l 4.333333 5.583333 l 4.166667 5.541667 l 4.125000 5.375000 l 3.958333 5.375000 l 3.958333 5.250000 l 3.791667 5.208333 l 3.791667 5.041667 l 3.625000 5.000000 l 3.583333 4.833333 l 3.416667 4.833333 l 3.375000 4.666667 l 2.875000 4.666667 l 2.833333 4.708333 l 2.833333 4.833333 l 3.041667 4.875000 l 3.041667 5.000000 l 3.208333 5.041667 l 3.208333 5.208333 l 3.375000 5.250000 l 3.416667 5.416667 l 3.583333 5.416667 l 3.583333 5.541667 l 3.750000 5.583333 l 3.750000 5.750000 l 3.916667 5.791667 l 3.958333 5.958333 l 4.125000 5.958333 l 4.125000 6.083333 l 4.291667 6.125000 l 4.291667 6.291667 l 4.458333 6.333333 l 4.500000 6.500000 l 4.666667 6.500000 l 4.666667 6.625000 l 4.833333 6.666667 l 4.833333 6.833333 l 5.000000 6.875000 l 5.041667 7.041667 l 5.208333 7.041667 l 5.250000 7.208333 l 5.541667 7.208333 l 5.583333 7.416667 l h f* Q}\special{pdf:econtent}\kern8.625000bp}}
\newcommand{\CometThreeSixNineD}{\hbox{\vrule width0pt height8.250000bp depth0pt\special{pdf:bcontent}\special{pdf:code q 0 g 8.416667 7.708333 m 8.083333 7.708333 l 8.083333 7.583333 l 7.916667 7.541667 l 7.916667 7.375000 l 7.750000 7.333333 l 7.708333 7.166667 l 7.541667 7.166667 l 7.541667 7.041667 l 7.375000 7.000000 l 7.375000 6.833333 l 7.208333 6.791667 l 7.166667 6.625000 l 6.833333 6.625000 l 6.833333 6.500000 l 6.666667 6.458333 l 6.666667 6.291667 l 6.458333 6.250000 l 6.458333 6.125000 l 6.291667 6.083333 l 6.250000 5.916667 l 6.125000 5.916667 l 6.125000 5.750000 l 5.958333 5.708333 l 5.916667 5.541667 l 5.750000 5.541667 l 5.750000 5.416667 l 5.583333 5.375000 l 5.583333 5.208333 l 5.416667 5.166667 l 5.375000 5.000000 l 5.208333 5.000000 l 5.208333 4.875000 l 5.041667 4.833333 l 5.041667 4.666667 l 4.875000 4.625000 l 4.833333 4.458333 l 4.500000 4.458333 l 4.500000 4.333333 l 4.333333 4.291667 l 4.333333 4.125000 l 4.125000 4.083333 l 4.125000 3.958333 l 3.958333 3.916667 l 3.916667 3.750000 l 3.583333 3.750000 l 3.541667 3.791667 l 3.541667 4.083333 l 3.750000 4.125000 l 3.791667 4.333333 l 4.083333 4.333333 l 4.125000 4.500000 l 4.291667 4.500000 l 4.291667 4.625000 l 4.458333 4.666667 l 4.458333 4.833333 l 4.625000 4.875000 l 4.666667 5.041667 l 4.833333 5.041667 l 4.833333 5.166667 l 5.000000 5.208333 l 5.000000 5.375000 l 5.166667 5.416667 l 5.208333 5.583333 l 5.375000 5.583333 l 5.375000 5.708333 l 5.541667 5.750000 l 5.541667 5.916667 l 5.708333 5.958333 l 5.750000 6.125000 l 5.916667 6.125000 l 5.958333 6.291667 l 6.250000 6.291667 l 6.250000 6.458333 l 6.416667 6.500000 l 6.458333 6.666667 l 6.625000 6.666667 l 6.625000 6.791667 l 6.791667 6.833333 l 6.791667 7.000000 l 6.958333 7.041667 l 7.000000 7.208333 l 7.166667 7.208333 l 7.166667 7.333333 l 7.333333 7.375000 l 7.375000 7.583333 l 7.500000 7.583333 l 7.500000 8.041667 l 6.833333 8.041667 l 6.833333 7.916667 l 6.666667 7.875000 l 6.625000 7.708333 l 6.291667 7.708333 l 6.291667 7.583333 l 6.125000 7.541667 l 6.125000 7.375000 l 5.958333 7.333333 l 5.958333 7.041667 l 5.750000 7.000000 l 5.750000 6.833333 l 5.583333 6.791667 l 5.583333 6.666667 l 5.416667 6.625000 l 5.375000 6.458333 l 5.041667 6.458333 l 5.041667 6.291667 l 4.875000 6.250000 l 4.833333 6.083333 l 4.666667 6.083333 l 4.666667 5.958333 l 4.500000 5.916667 l 4.500000 5.750000 l 4.333333 5.708333 l 4.291667 5.541667 l 4.125000 5.541667 l 4.125000 5.416667 l 3.958333 5.375000 l 3.958333 5.208333 l 3.791667 5.166667 l 3.750000 5.000000 l 3.583333 5.000000 l 3.583333 4.666667 l 3.541667 4.625000 l 3.250000 4.625000 l 3.208333 4.458333 l 3.041667 4.458333 l 3.000000 4.500000 l 3.000000 4.833333 l 3.208333 4.875000 l 3.208333 5.166667 l 3.375000 5.208333 l 3.375000 5.375000 l 3.541667 5.416667 l 3.583333 5.583333 l 3.916667 5.583333 l 3.916667 5.708333 l 4.083333 5.750000 l 4.083333 5.916667 l 4.291667 5.958333 l 4.291667 6.083333 l 4.458333 6.125000 l 4.500000 6.291667 l 4.625000 6.291667 l 4.625000 6.458333 l 4.833333 6.500000 l 4.833333 6.625000 l 5.000000 6.666667 l 5.000000 7.000000 l 5.166667 7.041667 l 5.208333 7.208333 l 5.375000 7.208333 l 5.375000 7.333333 l 5.541667 7.375000 l 5.541667 7.541667 l 5.708333 7.583333 l 5.750000 7.750000 l 5.916667 7.750000 l 5.916667 8.041667 l 5.750000 8.041667 l 5.708333 8.250000 l 5.208333 8.250000 l 5.166667 8.041667 l 4.875000 8.041667 l 4.833333 7.875000 l 4.666667 7.875000 l 4.666667 7.750000 l 4.500000 7.708333 l 4.458333 7.541667 l 4.333333 7.541667 l 4.333333 7.208333 l 4.125000 7.166667 l 4.125000 7.041667 l 3.958333 7.000000 l 3.958333 6.833333 l 3.791667 6.791667 l 3.750000 6.625000 l 3.583333 6.625000 l 3.583333 6.291667 l 3.416667 6.250000 l 3.416667 6.125000 l 3.250000 6.083333 l 3.208333 5.916667 l 3.041667 5.916667 l 3.041667 5.750000 l 2.875000 5.708333 l 2.875000 5.583333 l 2.708333 5.541667 l 2.666667 5.375000 l 2.500000 5.375000 l 2.500000 5.041667 l 2.333333 5.000000 l 2.291667 4.833333 l 2.166667 4.833333 l 2.166667 4.666667 l 1.958333 4.625000 l 1.916667 4.458333 l 1.416667 4.458333 l 1.375000 4.291667 l 1.083333 4.291667 l 1.083333 4.125000 l 0.916667 4.083333 l 0.875000 3.916667 l 0.708333 3.916667 l 0.708333 3.791667 l 0.541667 3.750000 l 0.500000 3.541667 l 0.375000 3.541667 l 0.375000 3.083333 l 0.166667 3.041667 l 0.166667 2.708333 l 0.000000 2.666667 l 0.000000 1.625000 l 0.166667 1.583333 l 0.166667 1.250000 l 0.375000 1.208333 l 0.375000 0.916667 l 0.541667 0.875000 l 0.541667 0.708333 l 0.666667 0.708333 l 0.708333 0.541667 l 1.041667 0.541667 l 1.083333 0.375000 l 1.375000 0.375000 l 1.416667 0.166667 l 2.125000 0.166667 l 2.166667 0.000000 l 2.666667 0.000000 l 2.708333 0.166667 l 3.000000 0.166667 l 3.041667 0.375000 l 3.375000 0.375000 l 3.375000 0.500000 l 3.541667 0.541667 l 3.583333 0.708333 l 3.750000 0.708333 l 3.750000 0.875000 l 3.916667 0.916667 l 3.958333 1.083333 l 4.083333 1.083333 l 4.083333 1.583333 l 4.291667 1.625000 l 4.291667 2.666667 l 4.083333 2.708333 l 4.083333 3.208333 l 4.291667 3.250000 l 4.291667 3.375000 l 4.458333 3.416667 l 4.500000 3.583333 l 4.833333 3.583333 l 4.833333 3.750000 l 5.000000 3.791667 l 5.000000 3.916667 l 5.166667 3.958333 l 5.208333 4.125000 l 5.541667 4.125000 l 5.541667 4.291667 l 5.708333 4.333333 l 5.750000 4.500000 l 5.916667 4.500000 l 5.958333 4.666667 l 6.250000 4.666667 l 6.250000 4.833333 l 6.416667 4.875000 l 6.458333 5.041667 l 6.625000 5.041667 l 6.625000 5.166667 l 6.791667 5.208333 l 6.833333 5.416667 l 7.166667 5.416667 l 7.166667 5.541667 l 7.333333 5.583333 l 7.375000 5.750000 l 7.500000 5.750000 l 7.500000 5.916667 l 7.708333 5.958333 l 7.708333 6.083333 l 7.875000 6.125000 l 7.916667 6.291667 l 8.041667 6.291667 l 8.083333 6.500000 l 8.416667 6.500000 l 8.416667 6.625000 l 8.583333 6.666667 l 8.583333 7.541667 l 8.458333 7.541667 l h 1.625000 3.791667 m 2.125000 3.791667 l 2.166667 3.583333 l 2.833333 3.583333 l 2.875000 3.416667 l 3.041667 3.375000 l 3.041667 3.250000 l 3.208333 3.250000 l 3.250000 3.083333 l 3.416667 3.041667 l 3.416667 2.333333 l 3.250000 2.291667 l 3.250000 1.791667 l 3.041667 1.750000 l 3.041667 1.458333 l 2.875000 1.416667 l 2.875000 1.250000 l 2.708333 1.208333 l 2.666667 1.041667 l 2.166667 1.041667 l 2.125000 0.875000 l 1.416667 0.875000 l 1.375000 1.041667 l 1.250000 1.041667 l 1.208333 1.208333 l 1.041667 1.250000 l 1.041667 1.416667 l 0.875000 1.458333 l 0.875000 1.750000 l 0.666667 1.791667 l 0.666667 2.833333 l 0.875000 2.875000 l 0.875000 3.041667 l 1.041667 3.083333 l 1.041667 3.375000 l 1.083333 3.416667 l 1.375000 3.416667 l 1.416667 3.583333 l 1.583333 3.583333 l h f* Q}\special{pdf:econtent}\kern8.583333bp}}

\par\indent Durch Kreismicrometer erhielt Herr Professor \textsc{Gauss} fünf Orte des Cometen;
\begin{center}\fontsize{9}{11}\selectfont\setlength{\tabcolsep}{4pt}\renewcommand{\arraystretch}{1.15}\begin{tabular}{c|c|c|c}
1813 & \shortstack{Mittlere Zeit\\in Göttingen} & AR. app. \CometThreeSixNineA & Decl. app. \CometThreeSixNineB\\\hline
April 7 & $13^{\mathrm u}\;12^{\mathrm m}\;2^{\mathrm s}$ & $271^\circ\;7'\;19''3$ & $5^\circ\;34'\;36''7$ N.\\
9 & $13\;35\;40$ & $270\;10\;33.5$ & $4\;11\;3.4$\\
11 & $13\;17\;43$ & $269\;1\;19.9$ & $2\;33\;0.7$\\
14 & $13\;7\;36$ & $266\;44\;5.5$ & $0\;33\;0.8$ S.\\
21 & $14\;23\;0$ & $256\;39\;19.3$ & $12\;57\;56.0$\\
\end{tabular}\end{center}
\Needspace{8\baselineskip}\par\indent Von \textsc{Harding} erhielten wir eine Meridianbeobachtung des Cometen
\begin{center}\fontsize{9}{11}\selectfont\setlength{\tabcolsep}{4pt}\renewcommand{\arraystretch}{1.15}\begin{tabular}{c|c|c|c}
1813 & \shortstack{Mittlere Zeit\\in Göttingen} & AR. app. \CometThreeSixNineC & Decl. app. \CometThreeSixNineD\\\hline
April 21 & $15^{\mathrm u}\;7^{\mathrm m}\;71^{\mathrm s}$ & $256^\circ\;34'\;19''6$ & $13^\circ\;2'\;26''5$ S.\\
\end{tabular}\end{center}

\Needspace{17\baselineskip}\SourceLeaf{370}\par\indent Aus den Beobachtungen vom 7. 9. 11. April berechnete Prof. \textsc{Gauss} und sein geschickter Schüler Hr. \textsc{Enke} folgende Elemente:

\begin{center}\fontsize{9}{11}\selectfont\setlength{\tabcolsep}{5pt}\renewcommand{\arraystretch}{1.15}\begin{tabular}{lcrr}
& & \textsc{Enke} & \textsc{Gauss}\\
Durchgang durchs Perihel 1813 & & Mai $19.7988$ & $20.3597$\\
Länge des Perihels & $.\;.\;.\;.\;.\;.$ & $196^\circ\;51'\;47''$ & $195^\circ\;24'\;57''$\\
Knoten & $.\;.\;.\;.\;.\;.\;.\;.\;.\;.\;.$ & $42\;29\;23$ & $42\;9\;26$\\
Neigung & $.\;.\;.\;.\;.\;.\;.\;.\;.\;.$ & $80\;49\;9$ & $80\;20\;11$\\
Logarithm. des kleinsten Abstandes &  & $0.08209$ & $0.07734$\\
\end{tabular}\end{center}

\Needspace{16\baselineskip}\par\indent Bald nachher erhielten wir von \textsc{Gauss} zweite, aus den Beobachtungen vom 7. 11. und 14. April berechnete Elemente:

\begin{center}\fontsize{10}{12}\selectfont\setlength{\tabcolsep}{2pt}\renewcommand{\arraystretch}{1.15}\begin{tabular}{l}
Durchgang durchs Perihelium 1813 Mai 15.93437\\Mittlere Zeit in Göttingen\\
\end{tabular}\end{center}\begin{center}\fontsize{10}{12}\selectfont\setlength{\tabcolsep}{2pt}\renewcommand{\arraystretch}{1.15}\begin{tabular}{lcr}
Länge des Perihels & $.\;.\;.\;.\;.\;.\;.\;.\;.\;.\;.\;.\;.$ & $196^\circ\;42'\;39''$\\
Logarithm. kleinster Abstand & $.\;.\;.\;.\;.\;.\;.\;.\;.$ & $0.08173$\\
Knoten & $.\;.\;.\;.\;.\;.\;.\;.\;.\;.\;.\;.\;.\;.\;.\;.\;.$ & $42^\circ\;29'\;12''$\\
Neigung & $.\;.\;.\;.\;.\;.\;.\;.\;.\;.\;.\;.\;.\;.\;.\;.\;.$ & $80\;46\;26$\\
\end{tabular}\end{center}

\noindent Bewegung rückläufig.

\par\indent Auch diese Elemente, schrieb uns letzterer bei deren Übersendung, bedürfen schon wieder einer Verbesserung, da sie am 21. April schon $9'$ differirten. Herr \textsc{Enke} ist jetzt mit dieser Verbesserung beschäftigt.
